% concatenated from article.tex
\documentclass[a4paper,10pt]{article}

\usepackage{authblk}

\usepackage[utf8]{inputenc} % allow utf-8 input
% \usepackage[frenchb]{babel}
% \usepackage[latin1]{inputenc}
\usepackage[T1]{fontenc}    % use 8-bit T1 fonts
\usepackage{indentfirst}
\usepackage{lmodern}

\usepackage{mathtools}
\usepackage{verbatim}
\usepackage{url}            % simple URL typesetting
\usepackage{booktabs}       % professional-quality tables
\usepackage{amsfonts}       % blackboard math symbols
\usepackage{amsmath}
\usepackage{amssymb}
\usepackage{amsthm}
\usepackage{microtype}      % microtypography
\usepackage[a4paper]{geometry} % reduce margins
\usepackage{enumitem}
\usepackage{xcolor}
\usepackage{mathrsfs}
\usepackage{calligra}

\usepackage{geometry}
%\geometry{margin=2.5cm}
\usepackage{fancyhdr}
\usepackage{listings}

\usepackage[linktocpage=true]{hyperref}
\hypersetup{
    colorlinks=true,
    linkcolor=blue,
    citecolor=violet,
    urlcolor=blue
}

\usepackage{graphicx}
\usepackage[makeroom]{cancel}
\usepackage{tikz}

\usepackage[inkscapeformat=png]{svg}
\usepackage[normalem]{ulem}
\usepackage{stmaryrd}



\sloppy

\theoremstyle{plain}
\newtheorem{theorem}{Theorem}[section]
\newtheorem{proposition}[theorem]{Proposition}
\newtheorem{corollary}[theorem]{Corollary}
\newtheorem{lemma}[theorem]{Lemma}

\theoremstyle{definition}
\newtheorem{definition}[theorem]{Definition}
\newtheorem{example}[theorem]{Example} 
\newtheorem{exercise}[theorem]{Exercise}
\newtheorem{problem}[theorem]{Problem}

\theoremstyle{remark}
\newtheorem{remark}[theorem]{Remark}

% \newenvironment{proof}{{\sc Proof:}}{~\hfill QED}

\DeclareMathOperator*{\essinf}{ess\,inf}
\DeclareMathOperator*{\esssup}{ess\,sup}
\DeclareMathOperator{\Gram}{Gram}
\DeclareMathOperator{\Span}{span}
\DeclareMathOperator{\Dim}{dim}
\DeclareMathOperator{\Ker}{Ker}
\DeclareMathOperator{\Gr}{Gr}      

\newcommand{\cI}{\mathcal{I}}
\newcommand{\Top}{\top}
\newcommand{\Bot}{\perp}
\newcommand{\vsim}{\simeq}
\newcommand{\cmp}{\circ}
\newcommand{\vh}{\star}
\newcommand{\vv}{\circ_{v}}
\newcommand{\cc}[2]{m'_{#1 #2}}
\newcommand{\Hom}[2]{\cI(S')(#1,\,#2)}           

\makeatletter
\def\keywords{\xdef\@thefnmark{}\@footnotetext}
\makeatother

\newcommand*\circled[1]{\tikz[baseline=(char.base)]{\node[shape=circle,draw,inner sep=2pt] (char) {#1};}}

%%%%%%%%%%%%%%%%%%%%%%%%%%      
%%%%%%%%%%%%%%%%%%%%%%%%%%
%%%%%%%%%%%%%%%%%%%%%%%%%%

\title{A general equivalence relation on the global class of morphisms of a category}
\author{Nizar El Idrissi}

\newcommand{\Addresses}{{% additional braces for segregating \footnotesize
  \bigskip
  \footnotesize

  \textbf{Nizar El Idrissi.}
  \par\nopagebreak Laboratoire : Equations aux dérivées partielles, Algèbre et Géométrie spectrales.
  \par\nopagebreak
  Département de mathématiques, faculté des sciences, université Ibn Tofail, 14000 Kénitra.\par\nopagebreak 
  \textit{E-mail address} : \texttt{nizar.elidrissi@uit.ac.ma}
  
%  \medskip
%  
%  \textbf{Hicham Zoubeir.}
%  \par\nopagebreak Laboratoire : Equations aux dérivées partielles, Algèbre et Géométrie spectrales.
%  \par\nopagebreak
%  Département de mathématiques, faculté des sciences, université Ibn Tofail, 14000 Kénitra.\par\nopagebreak 
%  \textit{E-mail address} : \texttt{hzoubeir2014@gmail.com}

}}

%%%%%%%%%%%%%%%%%%%%%
%%%%%%%%%%%%%%%%%%%%%
%%%%%%%%%%%%%%%%%%%%%


\begin{document}

\maketitle

\begin{abstract}
  We introduce a general equivalence relation on the global class of morphisms of a category that subsumes several classical notions of equivalence in mathematics. We show that group-action equivalence relations, pullback equivalence relations, and coarse analysis-operator-based equivalence relations on continuous Bessel families are special cases of this general definition.
\end{abstract}

\keywords{2020 \emph{Mathematics Subject Classification.} 18A99, 18B40, 46M15, 42C15 }
\keywords{\emph{Keywords and phrases.} category, 2-category, functor, equivalence relation, group-action, pullback, Bessel family}

\tableofcontents

\section{Introduction}

An equivalence relation on a set $X$ is a binary relation
$E \subseteq X \times X$ that is reflexive, symmetric, and transitive;
the equivalence class of a point $x \in X$ is denoted $[x]_E$,
and the quotient set $X/E$ collects all such classes.
Classical examples are ubiquitous: congruence modulo $n$ on $\mathbb{Z}$, the identification of Cauchy sequences of rationals used to construct
$\mathbb{R}$, isomorphism of groups or modules, and homotopy equivalence of topological spaces. Their systematic study in the Borel and analytic categories is developed at length in \cite{Kanoveui2008}. \\ \\
We can find many instances of such equivalence relations in functional analysis. Therein, a common way to define equivalence relations is to declare that two functions (or families) are equivalent if their respective inputs and outputs are interchangeable under a group action. \\ \\
Unitary equivalence of frames \cite{Christensen2016}, which dates back to the work of Han-Larson \cite{HanLarson2000}, is an equivalence relation of this sort. Of course, this equivalence relation can be further generalized to $\text{GL}(H)$ \cite{Balan1999} and its subgroups, such as $\text{L}^0(\Omega,\mathbb{S}^1) \ltimes \mathbb{U}(H)$ \cite{NovikovSevostyanova2002}. For example, Riesz bases are the sequences which are invertibly equivalent to an orthonormal basis. \\ \\
However, beyond group-action equivalence relations on function spaces, equivalence relations on morphism spaces naturally arise in categorical contexts where morphisms are compared via auxiliary categorical data rather than by automorphism groups. \\ \\
From a higher-categorical viewpoint, morphisms are naturally organized into 2-categories or bicategories where comparison data between morphisms is encoded as 2-morphisms. This perspective, developed systematically by many authors \cite{Benabou1967, Leinster2004}, provides a structural environment in which equivalence of morphisms is not necessarily induced by invertible transformations but may arise from more flexible coherence data. \\ \\
This suggests that there are many ways to define equivalence relations in mathematics, and that a unifying framework is desirable. \\ \\
\textbf{Main results}.
In this paper, we propose a general definition of equivalence relation on the global class of morphisms of a category, parametrized by functors and functions into a 2-category. This notion is reminiscent of Green's J relation \cite{CliffordPreston1961, Green1951, Howie1995} in semigroup theory. We recall that in semigroup theory, Green’s relations provide a paradigm where elements are classified through principal ideals and divisibility relations. The resulting equivalence relations depend on how elements factor through the ambient algebraic structure. \\
As an application, we show that group-action equivalence relations, set-theoretic pullback equivalence relations, and a coarse analysis-operator-based equivalence relation on the global class of continuous Bessel families are special cases of our general definition. \\ \\
\textbf{Plan of the article.} After this introduction, our paper is organized as follows. In section \ref{section-master}, we define an equivalence relation on the global class of morphisms of a category. In section \ref{section-examples}, we show that group-action equivalence relations, set-theoretic pullback equivalence relations, and a coarse analysis-operator-based equivalence relation on the global class of continuous Bessel families are special cases of our general definition. Section \ref{section-conclusion} contains the conclusion of this paper.

\section{A general categorical equivalence relation}
\label{section-master}

In this section, we propose an equivalence relation on the global class of morphisms of a category that unifies and extends several classical notions arising from group actions.

\begin{definition}
  \label{definition-master}
  Let $\mathcal{C}$ be a category, $\mathcal{D}$ be a 2-category. \\
  Let $\tau_1,\tau_2 : \mathcal{C} \to \mathcal{D}$ be two functors, and let
  \[ \sigma : \begin{cases} \mathsf{Ob(\mathcal{C})} &\to \mathsf{Ob(\mathcal{D})} \\ \mathsf{Hom(c_1,c_2)} &\to \mathsf{Hom(\sigma(c_1),\sigma(c_2))} \end{cases} \]
  be a function (not necessarily a functor) such that $\sigma, \tau_1$, and $\tau_2$ coincide on objects of $\mathcal{C}$. \\
  Further, let
  \[ m : A \to B \]
  and
  \[ \tilde{m} : \tilde{A} \to \tilde{B} \]
  be two morphisms in $\mathcal{C}$.  \\
  We say that $m$ and $\tilde{m}$ are equivalent (with respect to the data $(\mathcal{C},\mathcal{D},\sigma,\tau_1,\tau_2)$), or
  \[ m \simeq \tilde{m} \]
  if there exist 1-morphisms $u_1 : B \to \tilde{B}$, $u_2 : \tilde{A} \to A$, $v_1 : \tilde{B} \to B$, and $v_2 : A \to \tilde{A}$, and four 2-morphisms
  \[ \Phi : \tau_1(u_1) \circ \sigma(m) \circ \tau_2(u_2) \Rightarrow \sigma(\tilde{m}), \]
  \[ \tilde{\Phi} : \sigma(\tilde{m}) \Rightarrow \tau_1(u_1) \circ \sigma(m) \circ \tau_2(u_2), \]
  \[ \Psi : \tau_1(v_1) \circ \sigma(\tilde{m}) \circ \tau_2(v_2) \Rightarrow \sigma(m), \]
  and
  \[ \tilde{\Psi} : \sigma(m) \Rightarrow \tau_1(v_1) \circ \sigma(\tilde{m}) \circ \tau_2(v_2) . \]
\end{definition}

\begin{remark}
  \mbox{} \\
  \begin{itemize}
    \item The relation \(\simeq\) is defined on the class of all morphisms in the category \(\mathcal{C}\), with respect to the given data \((\mathcal{C}, \mathcal{D}, \sigma, \tau_1, \tau_2)\). \\
          Notice that the composites are well-defined since \(\sigma\), \(\tau_1\), and \(\tau_2\) coincide on objects, ensuring the domains and codomains align for the 1-morphisms in \(\mathcal{D}\).
    \item The definition of $  m \simeq \tilde{m}  $ does not require the $  1  $-morphisms $  u_1,u_2,v_1,v_2  $ (in $  \mathcal{C}  $) to be equivalences nor the $  2  $-morphisms $  \Phi,\tilde{\Phi},\Psi,\tilde{\Psi}  $ (in $  \mathcal{D}  $) to be isomorphisms. This is not a problem. Whenever one wishes to work in a setting where everything is invertible, it is always possible to replace $  \mathcal{D}  $ by a $  2  $-groupoid (a $  2  $-category in which every $  1  $-morphism is an equivalence and every $  2  $-morphism is an isomorphism), for example the $  2  $-categorical localisation of $  \mathcal{D}  $ obtained by formally inverting all $  1  $- and $  2  $-morphisms).
  \end{itemize}
\end{remark}

\begin{proposition}
  $\simeq$ is an equivalence relation.
\end{proposition}

\begin{proof}
  To prove \(\simeq\) is an equivalence relation, we show that it is reflexive, symmetric, and transitive. All constructions use the operations in the 2-category \(\mathcal{D}\): vertical composition \(\circ_v\) of 2-morphisms and whiskering (horizontal composition $\circ_h$) of a 2-morphism with 1-morphisms on the left and/or right.
  \begin{itemize}
    \item \textbf{Reflexivity.} For any morphism \(m: A \to B\), let's prove that \(m \simeq m\). \\
          Set \(\tilde{A} = A\), \(\tilde{B} = B\), \(\tilde{m} = m\). \\
          Choose \(u_1 = v_1 = \mathrm{id}_B: B \to B\), \(u_2 = v_2 = \mathrm{id}_A: A \to A\). \\
          Since \(\tau_1\) and \(\tau_2\) are functors, \(\tau_1(\mathrm{id}_B) = \mathrm{id}_{\tau_1(B)}\) and \(\tau_2(\mathrm{id}_A) = \mathrm{id}_{\tau_2(A)}\), and \(\sigma(A) = \tau_1(A) = \tau_2(A)\), \(\sigma(B) = \tau_1(B) = \tau_2(B)\).  \\
          Thus, \(\tau_1(u_1) \circ \sigma(m) \circ \tau_2(u_2) = \mathrm{id}_{\sigma(B)} \circ \sigma(m) \circ \mathrm{id}_{\sigma(A)} = \sigma(m)\). \\
          Similarly, \(\tau_1(v_1) \circ \sigma(m) \circ \tau_2(v_2) = \sigma(m)\). \\
          Set \(\Phi\), \(\tilde{\Phi}\), \(\Psi\), and \(\tilde{\Psi}\) to be the identity 2-morphism \(\mathrm{id}_{\sigma(m)}\) on the 1-morphism \(\sigma(m): \sigma(A) \to \sigma(B)\) in \(\mathcal{D}\).
          This satisfies the conditions.

    \item \textbf{Symmetry.} If \(m \simeq \tilde{m}\), let's prove that \(\tilde{m} \simeq m\). \\
          We are given \(u_1: B \to \tilde{B}\), \(u_2: \tilde{A} \to A\), \(v_1: \tilde{B} \to B\), \(v_2: A \to \tilde{A}\), and the 2-morphisms \(\Phi\), \(\tilde{\Phi}\), \(\Psi\), \(\tilde{\Psi}\). \\
          For \(\tilde{m} \simeq m\), set \(u_1': \tilde{B} \to B = v_1\), \(u_2': A \to \tilde{A} = v_2\), \(v_1': B \to \tilde{B} = u_1\), \(v_2': \tilde{A} \to A = u_2\). \\
          Then:
          \[ \Phi': \tau_1(u_1') \circ \sigma(\tilde{m}) \circ \tau_2(u_2') \Rightarrow \sigma(m) \text{ is }\Psi: \tau_1(v_1) \circ \sigma(\tilde{m}) \circ \tau_2(v_2) \Rightarrow \sigma(m)  \]
          \[ \tilde{\Phi}': \sigma(m) \Rightarrow \tau_1(u_1') \circ \sigma(\tilde{m}) \circ \tau_2(u_2') \text{ is } \tilde{\Psi}: \sigma(m) \Rightarrow \tau_1(v_1) \circ \sigma(\tilde{m}) \circ \tau_2(v_2) \]
          \[ \Psi': \tau_1(v_1') \circ \sigma(m) \circ \tau_2(v_2') \Rightarrow \sigma(\tilde{m}) \text{ is } \Phi: \tau_1(u_1) \circ \sigma(m) \circ \tau_2(u_2) \Rightarrow \sigma(\tilde{m}) \]
          \[ \tilde{\Psi}': \sigma(\tilde{m}) \Rightarrow \tau_1(v_1') \circ \sigma(m) \circ \tau_2(v_2') \text{ is } \tilde{\Phi}: \sigma(\tilde{m}) \Rightarrow \tau_1(u_1) \circ \sigma(m) \circ \tau_2(u_2). \]
          This satisfies the conditions by relabeling.

    \item \textbf{Transitivity.} If \(m \simeq \overline{m}\) and \(\overline{m} \simeq \overline{\overline{m}}\), let's prove that \(m \simeq \overline{\overline{m}}\). \\
          Let \(m: A \to B\), \(\overline{m}: C \to D\), \(\overline{\overline{m}}: E \to F\). \\
          Since \(m \simeq  \overline{m}\), there exist \(u_1: B \to D\), \(u_2: C \to A\), \(\Phi: \tau_1(u_1) \circ \sigma(m) \circ \tau_2(u_2) \Rightarrow \sigma(\overline{m})\), \(\tilde{\Phi}: \sigma(g) \Rightarrow \tau_1(u_1) \circ \sigma(m) \circ \tau_2(u_2)\), \(v_1: D \to B\), \(v_2: A \to C\), \(\Psi: \tau_1(v_1) \circ \sigma(\overline{m}) \circ \tau_2(v_2) \Rightarrow \sigma(m)\), \(\tilde{\Psi}: \sigma(m) \Rightarrow \tau_1(v_1) \circ \sigma(\overline{m}) \circ \tau_2(v_2)\). \\
          Since \(\overline{m} \simeq \overline{\overline{m}}\), there exist \(u_1': D \to F\), \(u_2': E \to C\), \(\Phi': \tau_1(u_1') \circ \sigma(\overline{m}) \circ \tau_2(u_2') \Rightarrow \sigma(\overline{\overline{m}})\), \(\tilde{\Phi}': \sigma(\overline{\overline{m}}) \Rightarrow \tau_1(u_1') \circ \sigma(\overline{m}) \circ \tau_2(u_2')\), \(v_1': F \to D\), \(v_2': C \to E\), \(\Psi': \tau_1(v_1') \circ \sigma(\overline{\overline{m}}) \circ \tau_2(v_2') \Rightarrow \sigma(\overline{m})\), \(\tilde{\Psi}': \sigma(g) \Rightarrow \tau_1(v_1') \circ \sigma(\overline{\overline{m}}) \circ \tau_2(v_2')\). \\
          Now, for \(m \simeq \overline{\overline{m}} \), set \(u_1'': B \to F = u_1' \circ u_1\), \(u_2'': E \to A = u_2 \circ u_2'\), \(v_1'': F \to B = v_1 \circ v_1'\), \(v_2'': A \to E = v_2' \circ v_2\). \\
          Since \(\tau_1\) and \(\tau_2\) are functors, \(\tau_1(u_1'' ) = \tau_1(u_1') \circ \tau_1(u_1)\), \(\tau_2(u_2'') = \tau_2(u_2) \circ \tau_2(u_2')\), \(\tau_1(v_1'') = \tau_1(v_1) \circ \tau_1(v_1')\), \(\tau_2(v_2'') = \tau_2(v_2') \circ \tau_2(v_2)\). \\
          Construct the 2-morphisms:
          \begin{itemize}
            \item \(\Phi'': \tau_1(u_1'') \circ \sigma(m) \circ \tau_2(u_2'') \Rightarrow \sigma(\overline{\overline{m}})\) is the vertical composition
                  \[ \Phi' \circ_v \left( 1_{\tau_1(u_1')} \circ_h \Phi \circ_h 1_{\tau_2(u_2')} \right), \]
                  where $1_{\tau_1(u_1')} \circ_h \Phi \circ_h 1_{\tau_2(u_2')}$ is \(\Phi\) whiskered on the left by \(\tau_1(u_1')\) and on the right by \(\tau_2(u_2')\).
            \item \(\tilde{\Phi}'': \sigma(\overline{\overline{m}}) \Rightarrow \tau_1(u_1'') \circ \sigma(m) \circ \tau_2(u_2'')\) is the vertical composition
                  \[ \left( 1_{\tau_1(u_1')} \circ_h \tilde{\Phi} \circ_h 1_{\tau_2(u_2')} \right) \circ_v \tilde{\Phi}', \]
                  where $1_{\tau_1(u_1')} \circ_h \tilde{\Phi} \circ_h 1_{\tau_2(u_2')}$ is \(\tilde{\Phi}\) whiskered on the left by \(\tau_1(u_1')\) and on the right by \(\tau_2(u_2')\).
            \item \(\Psi'': \tau_1(v_1'') \circ \sigma(\overline{\overline{m}}) \circ \tau_2(v_2'') \Rightarrow \sigma(m)\) is the vertical composition
                  \[ \Psi \circ_v \left( 1_{\tau_1(v_1)} \circ_h \Psi' \circ_h 1_{\tau_2(v_2)} \right), \]
                  where $1_{\tau_1(v_1)} \circ_h \Psi' \circ_h 1_{\tau_2(v_2)}$ is \(\Psi'\) whiskered on the left by \(\tau_1(v_1)\) and on the right by \(\tau_2(v_2)\).
            \item \(\tilde{\Psi}'': \sigma(m) \Rightarrow \tau_1(v_1'') \circ \sigma(\overline{\overline{m}}) \circ \tau_2(v_2'')\) is the vertical composition
                  \[ \left( 1_{\tau_1(v_1)} \circ_h \tilde{\Psi}' \circ_h 1_{\tau_2(v_2)} \right) \circ_v \tilde{\Psi}, \] where $1_{\tau_1(v_1)} \circ_h \tilde{\Psi}' \circ_h 1_{\tau_2(v_2)}$ is \(\tilde{\Psi}'\) whiskered on the left by \(\tau_1(v_1)\) and on the right by \(\tau_2(v_2)\).
          \end{itemize}
          This satisfies the conditions.
  \end{itemize}
\end{proof}





\section{Examples}
\label{section-examples}

\subsection{Group-action equivalence relation}

\begin{definition}[Group-action equivalence]
  \label{definition-group-equivalence}
  Let $G$ be a group acting on a set $E$.
  Define the \textbf{group-action equivalence relation} $\simeq_G$ on $E$ by
  \[
    f :\simeq_G \tilde{f} \quad \Longleftrightarrow \quad \exists g \in G \text{ such that } \tilde{f} = g. f
  \]
\end{definition}

\mbox{} \\
Now, let's show that the equivalence relation of definition \ref{definition-group-equivalence} fits into the general framework of definition \ref{definition-master}. \\
Consider the identity functors $\sigma = \tau_1 = \tau_2 = \mathsf{Id}$ on $\mathsf{\textbf{B} \left( (E \sslash G)^{\otimes *} \right) }$, where $\mathsf{\textbf{B}\left( (E \sslash G)^{\otimes*} \right) }$ is the 2-category delooping of the free strict monoidal category generated by the action groupoid $\mathsf{E \sslash G}$.

\begin{proposition}
  For all $f \in E$, $f$ can be seen as a morphism in $\mathsf{\textbf{B} \left( (E \sslash G)^{\otimes *} \right) }$. Under this identification, $f \simeq_G \tilde{f}$ as in definition \ref{definition-group-equivalence} if and only if $f \simeq \tilde{f}$ with respect to the data $(\mathsf{\textbf{B} \left( (E \sslash G)^{\otimes *} \right) },\mathsf{\textbf{B} \left( (E \sslash G)^{\otimes *} \right) },\mathsf{Id},\mathsf{Id},\mathsf{Id})$, as in definition \ref{definition-master}.
\end{proposition}


\begin{proof}
  This follows from length comparison: for all $f, \tilde{f} \in E$, for all chains $u_1,u_2 \in E^*$ and for all $g \in G$, we have $u_1 f u_2 = g.\tilde{f}$ if and only if $u_1$ and $u_2$ are the empty chains and $f = g.\tilde{f}$.
\end{proof}


\subsection{Pullback equivalence relation}

\begin{definition}
  \label{definition-pullback-equivalence-relation}
Let $S$ and $S'$ be sets, where $S'$ is equipped with an equivalence relation $\simeq$. Let $f : S \to S'$ be a function. Define the \textbf{pullback equivalence relation} $\simeq_f$ on $S$ by 
\[
  s \simeq_f \tilde{s} \quad \Longleftrightarrow \quad f(s) \simeq f(\tilde{s}).
\]
If $\simeq$ is seen as a subset of $S' \times S'$, then $\simeq_f$ is the subset of $S \times S$ defined by $\simeq_f = (f^{\triangleq})^{-1}(\simeq)$, where $f^{\triangleq} : S \times S \to S' \times S'$ is the function defined by $f^{\triangleq}(s,\tilde{s}) = (f(s), f(\tilde{s}))$.
\end{definition}

\begin{definition}
  Let $S'$ be a set equipped with an equivalence relation $\vsim$. \\
  Define the
  \emph{truth-value labeling} $\tau : S' \times S' \to \{\Top, \Bot\}$ by
  \[
    \tau(s', \tilde{s}') \;:=\;
    \begin{cases}
      \Top & \text{if } s' \vsim \tilde{s}', \\
      \Bot & \text{otherwise.}
    \end{cases}
  \]
  We construct a 2-category $\mathrm{Ind}(S')$, called the
  \textbf{indiscrete 2-category on $(S',\vsim)$}, as follows.
  \begin{itemize}
    \item \textbf{0-cells.} The elements of $S'$.
    \item \textbf{1-cells.} For each pair $s', \tilde{s}' \in S'$, a unique
          1-cell $\cc{s'}{\tilde{s}'} : s' \to \tilde{s}'$.
    \item \textbf{2-cells.} For each pair $s', \tilde{s}' \in S'$, the unique
          2-cell
          \[
            \Phi_{s'\tilde{s}'} \;:\; \cc{s'}{\tilde{s}'} \Rightarrow
            \cc{s'}{\tilde{s}'}
          \]
          is labeled by $\tau(s', \tilde{s}')$. We write this 2-cell as $\Top$ when
          $s' \vsim \tilde{s}'$, and as $\Bot$ when $s' \not\vsim \tilde{s}'$.
    \item \textbf{Composition of 1-cells.}
          $\cc{\tilde{s}'}{s''} \cmp \cc{s'}{\tilde{s}'} := \cc{s'}{s''}$.
    \item \textbf{Identity 1-cells.}
          $\mathrm{id}_{s'} := \cc{s'}{s'}$.

    \item \textbf{Identity 2-cells.} Define $\mathrm{id}_{\cc{s'}{\tilde{s}'}} :=
            \tau(s', \tilde{s}')$. This is forced, since there is only one 2-cell
          available.

    \item \textbf{Vertical composition of 2-cells.} 
          
    Define the vertical composition by
          \[
            \tau(s', \tilde{s}') \vv \tau(s', \tilde{s}') \;:=\; \tau(s', \tilde{s}').
          \]
          This is well-defined because all these 2-cells are equal to the the unique 2-cell in that hom-category.

          \begin{itemize}

            \item \textbf{Associativity.} For any triple of 2-cells (each equal to
          $\tau(s', \tilde{s}')$),
          \[
            \bigl(\tau(s',\tilde{s}') \vv \tau(s',\tilde{s}')\bigr) \vv \tau(s',\tilde{s}')
            \;=\; \tau(s',\tilde{s}')
            \;=\; \tau(s',\tilde{s}') \vv \bigl(\tau(s',\tilde{s}') \vv
            \tau(s',\tilde{s}')\bigr). \;
          \]

          \item \textbf{Unit laws.} Both left and right unit laws read
          \[
            \tau(s', \tilde{s}') \vv \tau(s', \tilde{s}') \;=\; \tau(s', \tilde{s}'),
          \]
          which holds by definition.

          \end{itemize}
        
      \item \textbf{Horizontal composition of 2-cells.}

          Define the horizontal composition by
          \[
            \tau(\tilde{s}', s'') \vh \tau(s', \tilde{s}') := \tau(s', s'').
          \]

          \begin{itemize}
            \item \textbf{Associativity.}
                  For any triple of composable 1-cells
                  $s' \to \tilde{s}' \to \tilde{\tilde{s}}' \to \tilde{\tilde{\tilde{s}}}'$, the associativity constraint reads
                  \[
                    (\tau(\tilde{\tilde{s}}', \tilde{\tilde{\tilde{s}}}') \vh \tau(\tilde{s}', \tilde{\tilde{s}}')) \vh \tau(s', \tilde{s}')
                    = \tau(s', \tilde{\tilde{\tilde{s}}}') = \tau(\tilde{\tilde{s}}', \tilde{\tilde{\tilde{s}}}') \vh (\tau(\tilde{s}', \tilde{\tilde{s}}') \vh
                    \tau(s', \tilde{s}')).
                  \]
            
            \item \textbf{Unit laws.}
                  $\mathrm{id}_{\cc{\tilde{s}'}{\tilde{\tilde{s}}'}} \vh \mathrm{id}_{\cc{s'}{\tilde{s}'}}
                    = \tau(\tilde{s}', \tilde{\tilde{s}}') \vh \tau(s', \tilde{s}')
                    = \tau(s', \tilde{\tilde{s}}')
                    = \mathrm{id}_{\cc{s'}{\tilde{\tilde{s}}'}}$.
          \end{itemize}


    \item \textbf{Interchange law.}

          For any composable configuration
          \[
            s' \xrightarrow{\;\cc{s'}{\tilde{s}'}\;} \tilde{s}'
            \xrightarrow{\;\cc{\tilde{s}'}{\tilde{\tilde{s}}'}\;} \tilde{\tilde{s}}'
          \]
          and vertically composable 2-cells $\Phi_1, \Phi_2$ above the first 1-cell and
          $\Psi_1, \Psi_2$ above the second, the interchange law states
          \[
            (\Psi_2 \vv \Psi_1) \vh (\Phi_2 \vv \Phi_1)
            \;=\;
            (\Psi_2 \vh \Phi_2) \vv (\Psi_1 \vh \Phi_1).
          \]
          Both sides equal the unique 2-cell $\tau(s', \tilde{\tilde{s}}')$ in $\mathrm{Ind}(S')(s',\tilde{\tilde{s}}')$.

       \item \textbf{Associativity constraint.}

          For any triple of composable 1-cells
          $s' \to \tilde{s}' \to \tilde{\tilde{s}}' \to \tilde{\tilde{\tilde{s}}}'$, the associator
          \[
            \alpha_{s',\tilde{s}',\tilde{\tilde{s}}',\tilde{\tilde{\tilde{s}}}'}
            \;:\;
            \bigl(\cc{\tilde{s}'}{\tilde{\tilde{s}}'} \cmp \cc{s'}{\tilde{s}'}\bigr) \cmp \cc{\tilde{\tilde{s}}'}{ \tilde{\tilde{\tilde{s}}}'}
            \;\Rightarrow\;
            \cc{\tilde{s}'}{\tilde{\tilde{s}}'} \cmp \bigl(\cc{s'}{\tilde{s}'} \cmp \cc{\tilde{\tilde{s}}'}{\tilde{\tilde{\tilde{s}}}'}\bigr)
          \]
          is a 2-cell in $\mathrm{Ind}(S')(s',\tilde{\tilde{\tilde{s}}}')$, hence is uniquely determined as $\tau(s', \tilde{\tilde{\tilde{s}}}')$.
          The pentagon coherence condition reduces to
          \[
            \tau(s', \tilde{\tilde{\tilde{\tilde{s}}}}') \;=\; \tau(s', \tilde{\tilde{\tilde{\tilde{s}}}}'),
          \]
          which holds trivially.

    \item \textbf{Unit constraints.}

          The left unitor $\lambda_{s'\tilde{s}'} : \cc{s'}{s'} \cmp \cc{s'}{\tilde{s}'}
            \Rightarrow \cc{s'}{\tilde{s}'}$ and the right unitor
          $\rho_{s'\tilde{s}'} : \cc{s'}{\tilde{s}'} \cmp \cc{\tilde{s}'}{\tilde{s}'}
            \Rightarrow \cc{s'}{\tilde{s}'}$ are both 2-cells in $\mathrm{Ind}(S')(s',\tilde{s}')$,
          hence each equals $\tau(s', \tilde{s}')$. The triangle coherence condition
          reads
          \[
            \tau(s', \tilde{\tilde{s}}') \;=\; \tau(s', \tilde{\tilde{s}}'),
          \]
          which again holds trivially.

  \end{itemize}

    Since all axioms are satisfied, $\mathrm{Ind}(S')$ is a (strict) 2-category.
\end{definition}

\begin{remark}[Role of the equivalence relation]
  The labeling $\tau$ stratifies the 2-cells into two classes:
  \begin{itemize}
    \item \textbf{Diagonal} ($s' \vsim \tilde{s}'$): the 2-cell $\Top$ plays the
          rôle of a ``true'' identity, reflecting that $s'$ and $\tilde{s}'$ are
          equivalent objects.
    \item \textbf{Off-diagonal} ($s' \not\vsim \tilde{s}'$): the 2-cell $\Bot$
          is a ``false'' or null morphism, signaling a transition between
          inequivalent classes.
  \end{itemize}
\end{remark}

Now, let $S$ and $S'$ be sets, where $S'$ is equipped with an equivalence relation $\simeq$. \\
Let $f : S \to S'$ be a function. Define the functor
\[
  \sigma \;=\; \tau_1 \;=\; \tau_2 \;:\; \mathrm{Ind}(S) \;\longrightarrow\; \mathrm{Ind}(S')
\]
by $\sigma(a) := f(a)$ on 0-cells and
$\sigma\!\left(m_{ab}\right) := m'_{f(a)f(b)}$ on 1-cells,
extended to 2-cells by $\sigma(\tau(a,b)) := \tau'(f(a),f(b))$.

\begin{proposition}
  \label{proposition-pullback-equivalence-relation-is-the-same-as-the-main-equivalence-relation}
  Let $s, \tilde{s} \in S$. Then
  \begin{align*}
    s \;\simeq_f\; \tilde{s} \text{ in } S \text{ (in the sense of definition \ref{definition-pullback-equivalence-relation})} &\iff  m_{s\tilde{s}} \;\simeq\; m_{s\tilde{s}}
    \quad \text{ with respect to the data } \\
    &(\mathrm{Ind}(S), \mathrm{Ind}(S'), \mathsf{\sigma}, \mathsf{\sigma}, \mathsf{\sigma}) \text{ (in the sense of definition \ref{definition-master})},
  \end{align*}
where $m_{s\tilde{s}} : s \to \tilde{s}$ is the unique
  1-cell in $\mathrm{Ind}(S)$. 
\end{proposition}

\begin{proof}

  \medskip
  \noindent\textbf{Reduction of the data.}

  \medskip
  We apply definition \ref{definition-master} with $m = \tilde{m} = m_{s\tilde{s}}$,
  $A = \tilde{A} = s$, and $B = \tilde{B} = \tilde{s}$.

  \smallskip
  \noindent\emph{Choice of auxiliary 1-cells.}
  The 1-cells $u_1 : \tilde{s} \to \tilde{s}$, $u_2 : s \to s$,
  $v_1 : \tilde{s} \to \tilde{s}$, $v_2 : s \to s$ required by
  definition~\ref{definition-master} are \emph{uniquely determined} in $\mathrm{Ind}(S)$:
  \[
    u_1 = m_{\tilde{s}\tilde{s}}, \quad
    u_2 = m_{ss}, \quad
    v_1 = m_{\tilde{s}\tilde{s}}, \quad
    v_2 = m_{ss}.
  \]
  There is no freedom of choice; any pair of objects in $\mathrm{Ind}(S)$ admits
  exactly one 1-cell between them.

  \smallskip
  \noindent\emph{Computation of the composite 1-cells in $\mathrm{Ind}(S')$.}
  Applying $\sigma = \tau_1 = \tau_2$ gives
  \[
    \sigma(u_1) = m'_{f(\tilde{s})f(\tilde{s})}, \quad
    \sigma(u_2) = m'_{f(s)f(s)}, \quad
    \sigma(m_{s\tilde{s}}) = m'_{f(s)f(\tilde{s})}.
  \]
  Since $\mathrm{Ind}(S')$ is indiscrete, the composite
  \[
    \sigma(u_1) \circ \sigma(m_{s\tilde{s}}) \circ \sigma(u_2)
    \;=\; m'_{f(\tilde{s})f(\tilde{s})} \circ
    m'_{f(s)f(\tilde{s})} \circ
    m'_{f(s)f(s)}
    \;=\; m'_{f(s)f(\tilde{s})}
    \;=\; \sigma(m_{s\tilde{s}}).
  \]
  By the same computation with $v_1, v_2$ (which equal $u_1, u_2$), the
  composites for $\Psi$ and $\widetilde{\Psi}$ likewise both equal
  $m'_{f(s)f(\tilde{s})}$.

  \smallskip
  \noindent\emph{The required 2-morphisms.}
  All four 2-morphisms $\Phi, \widetilde{\Phi}, \Psi, \widetilde{\Psi}$
  reduce to endomorphisms of the single 1-cell
  $m'_{f(s)f(\tilde{s})}$:
  \begin{align}
    \Phi,\,\widetilde{\Phi},\,\Psi,\,\widetilde{\Psi}
    \;:\;
    m'_{f(s)f(\tilde{s})} \Rightarrow m'_{f(s)f(\tilde{s})}
    \quad \text{in } \mathrm{Ind}(S').
    \label{eq:2cells}
  \end{align}
  The unique 2-cell in $\mathrm{Ind}(S')(f(s),\,f(\tilde{s}))$ is labeled
  $\tau'(f(s), f(\tilde{s})) \in \{\top, \bot\}$. By our convention, this
  constitutes a genuine 2-morphism if and only if
  $\tau'(f(s), f(\tilde{s})) = \top$, equivalently $f(s) \simeq f(\tilde{s})$
  in $S'$.

  \begin{itemize}
  
  \item \textbf{Proof of $(\Leftarrow)$: $f(s) \simeq f(\tilde{s})$
    implies $m_{s\tilde{s}} \simeq m_{s\tilde{s}}$.} \\
  Suppose $f(s) \simeq f(\tilde{s})$ in $S'$, so $\tau'(f(s),f(\tilde{s}))
    = \top$. Choose the auxiliary 1-cells as above. By~\eqref{eq:2cells}, each
  of the four required 2-morphisms is the unique 2-cell
  \[
    \Phi \;=\; \widetilde{\Phi} \;=\; \Psi \;=\; \widetilde{\Psi}
    \;:=\; \top
    \;\in\; \mathrm{Ind}(S')(f(s),\,f(\tilde{s})).
  \]
  Since $\tau'(f(s), f(\tilde{s})) = \top$, this is a genuine 2-morphism
  in $\mathrm{Ind}(S')$. Hence all four 2-morphisms of definition~\ref{definition-master}
  exist, and $m_{s\tilde{s}} \simeq m_{s\tilde{s}}$.

 \item \textbf{Proof of $(\Rightarrow)$: $m_{s\tilde{s}} \simeq m_{s\tilde{s}}$ implies
    $f(s) \simeq f(\tilde{s})$.}

  \medskip
  Suppose $m_{s\tilde{s}} \simeq m_{s\tilde{s}}$, so there exist 1-cells and 2-morphisms as in
  definition~\ref{definition-master}. As shown in the reduction step, the auxiliary
  1-cells are uniquely forced, and the four 2-morphisms all take the
  form~\eqref{eq:2cells}. In particular, $\Phi$ is a genuine 2-morphism in
  $\mathrm{Ind}(S')(f(s),\,f(\tilde{s}))$. The existence of any genuine 2-morphism in
  that hom-category requires $\tau'(f(s), f(\tilde{s})) = \top$, i.e.\
  $f(s) \simeq f(\tilde{s})$ in $S'$.
  \end{itemize}
  Both implications having been established, the proof is complete.
\end{proof}

\begin{remark}[Uniqueness forces the equivalence]
  The key mechanism is a two-fold uniqueness:
  \begin{enumerate}
    \item \emph{Uniqueness of 1-cells.} The indiscrete structure of
          $\mathrm{Ind}(S)$ forces $u_1, u_2, v_1, v_2$ to their only possible values,
          eliminating any flexibility in definition~\ref{definition-master}.
    \item \emph{Uniqueness of 2-cells.} Each hom-category of $\mathrm{Ind}(S')$
          contains at most one genuine 2-morphism. Whether that 2-morphism
          exists is controlled entirely by the truth value
          $\tau'(f(s), f(\tilde{s}))$, which in turn encodes the equivalence
          relation on $S'$.
  \end{enumerate}
  Together, these two uniqueness principles collapse the general
  definition~\ref{definition-master} into a single binary condition:
  $f(s) \simeq f(\tilde{s})$.
\end{remark}


\subsection{Coarse analysis-operator-based equivalence relation on the global class of continuous Bessel families}

\begin{definition}[Operatorial equivalence in varying $\mathsf{\mathcal{B}}(\Omega,H)$]
  \label{definition-operatorial-equivalence}
  For any two Hilbert spaces $H_1$ and $H_2$ over $\mathbb{F}$ ($\mathbb{R}$ or $\mathbb{C}$), and for any operator $T : H_1 \to H_2$, define
  \[ \rho(T) := \lVert T(\cdot) \rVert_{H_2} : H_1 \to \mathbb{R}_+. \]
  Let $H$ be a Hilbert space over $\mathbb{F}$ ($\mathbb{R}$ or $\mathbb{C}$) and let $(\Omega,\Sigma,\mu)$ be a measure space.\\
  For $f \in \mathsf{\mathcal{B}}(\Omega,H)$, let $T_f : H \to L^2(\Omega)$ be the analysis operator
  \[
    T_f(x) = (\langle x, f(\omega) \rangle)_{\omega \in \Omega},
  \]
  and define
  \[
    \boldsymbol{\rho}(f) := \rho(T_f) = \|T_f(\cdot)\|_{L_2(\Omega)} : H \to \mathbb{R}_+.
  \]
  Let $\mathsf{Hilb_{\mathbb{F}}^{\text{inj}}}$, $\mathsf{Hilb_{\mathbb{F}}^{\text{inj,cl}}}$, and $\mathsf{Hilb_{\mathbb{F}}^{\text{iso}}}$ be respectively the categories of Hilbert spaces with bounded injective operators, bounded injective operators with closed range, and isometries as morphisms. \\
  Define the \textbf{coarse analysis-operator-based equivalence relation} on the global class of continuous Bessel families as follows: two functions $f \in \mathsf{\mathcal{B}}(\Omega,H)$ and $\tilde{f} \in \mathsf{\mathcal{B}}(\tilde{\Omega},\tilde{H})$ are defined to be equivalent (with respect to the category of Hilbert spaces that has been chosen), denoted $f \simeq \tilde{f}$, if there exist morphisms $u : H \to \tilde{H}$ and $\tilde{u} : \tilde{H} \to H$ such that
  \[
    \boldsymbol{\rho}(u \circ f) \asymp \boldsymbol{\rho}(\tilde{f}) \quad \text{and} \quad \boldsymbol{\rho}(\tilde{u} \circ \tilde{f}) \asymp \boldsymbol{\rho}(f),
  \]
  or equivalently,
  \[
    \rho(T_f \circ u^*) \asymp \rho(T_{\tilde{f}}) \quad \text{and} \quad \rho(T_{\tilde{f}} \circ \tilde{u}^*) \asymp \rho(T_f),
  \]
  Here, $A \asymp B$ (with $A,B : S \to \mathbb{R}_+$ and $S$ an arbitrary set) means that there exist two uniform constants $K_1,K_2>0$ such that
  \[ \forall x \in S : \quad \quad K_1 A(x) \leq B(x) \leq K_2 A(x). \]
\end{definition}

\begin{remark}
  \mbox{} \\
  \begin{itemize}
    \item This equivalence relation ignores pointwise structure (we have heterogeneous indexing sets) and focuses exclusively on global function behavior encoded by the analysis operator. Moreover, this formulation allows for different and not necessarily isomorphic function codomains.  Therefore, it is coarser than invertible or unitary equivalence.
    \item For this equivalence relation, frames $f \in \mathsf{\mathcal{F}}(\Omega,H)$ are the functions that are equivalent to orthonormal bases $\tilde{f} : \operatorname{dim}(H) \to H$ (where $\operatorname{dim}(H)$ is the cardinality of $H$). \\
          Indeed, if $f \simeq \tilde{f}$ and $\tilde{f}$ is an orthonormal basis, then, in particular,
          \[ (\exists K_1,K_2 > 0)(\forall y \in \tilde{H}) : K_1  \lVert y \rVert^2 \leq \lVert T_f \circ u^*(y) \rVert^2 \leq K_2 \lVert y \rVert^2. \]
          This implies that $u \circ T_f^*$ is surjective, so $u$ is surjective, but since it is already injective, it is an isomorphism. \\
          Therefore $T_f^*$ is surjective, which means that $f \in \mathsf{\mathcal{F}}(\Omega,H)$.
  \end{itemize}
\end{remark}

\begin{proposition}
  \mbox{} \\
  \begin{itemize}
    \item For any measure space $(\Omega ,\Sigma,\mu)$, Hilbert spaces $H$ and $\tilde{H}$, $f \in \mathsf{\mathcal{B}}(\Omega,H)$, and $v : H \to \tilde{H}$, we have $T_f \circ \alpha = T_{\alpha^* \circ f}$. \\
          Consequently, we have
          \[ \boldsymbol{\rho}(f) \circ \alpha = \rho(T_f) \circ \alpha = \|(T_f \circ \alpha)(\cdot) \|_{L_2(\Omega)} = \|T_{\alpha^* \circ f}(\cdot) \|_{L_2(\Omega)} =\rho(T_{\alpha^* \circ f}) = \boldsymbol{\rho}(\alpha^* \circ f). \]
    \item $\simeq$ is an equivalence relation.
  \end{itemize}
\end{proposition}

\begin{proof}
  \mbox{} \\
  \begin{itemize}
    \item For any measure space $(\Omega ,\Sigma,\mu)$, Hilbert spaces $H$ and $\tilde{H}$, $f \in \mathsf{\mathcal{B}}(\Omega,H)$, $v : H \to \tilde{H}$, and $z \in \tilde{H}$, we have
          \[ (T_f\circ \alpha)(z)=T_f( \alpha z)=(\langle \alpha z,f(\omega)\rangle)_{\omega\in\Omega}=(\langle z, \alpha^* f(\omega)\rangle)_{\omega\in\Omega}=T_{\alpha^* \circ f}(z). \]
          The two formulations are therefore identical.


    \item \textbf{Proof that the relation is an equivalence relation.} \\
          The relation $\simeq$ is defined on the collection of all $f \in \mathsf{\mathcal{B}}(\Omega,H)$ (over all possible measure spaces $(\Omega,\Sigma,\mu)$ and Hilbert spaces $H$), where $f\simeq\tilde{f}$ if there exist bounded linear operators $u:H\to\tilde{H}$ and $\tilde{u}:\tilde{H}\to H$ such that $\boldsymbol{\rho}(u\circ f)\asymp\boldsymbol{\rho}(\tilde{f})$ and $\boldsymbol{\rho}(\tilde{u}\circ\tilde{f})\asymp\boldsymbol{\rho}(f)$. \\
          We verify reflexivity, symmetry, and transitivity.
          \begin{itemize}
            \item \textbf{Reflexivity.}  Fix $f\in \mathsf{\mathcal{B}}(\Omega,H)$. Take $u=\mathrm{Id}_H:H\to H$ and $\tilde{u}=\mathrm{Id}_H:H\to H$. Then $u\circ f=f$, so $\boldsymbol{\rho}(u\circ r)=\boldsymbol{\rho}(f)$, and similarly $\boldsymbol{\rho}(\tilde{u}\circ r)=\boldsymbol{\rho}(r)$. Thus, $\boldsymbol{\rho}(r)\asymp\boldsymbol{\rho}(r)$ holds with constants $K_1=K_2=1$ for all $x\in H$. Hence, $f\simeq r$.
            \item \textbf{Symmetry.}  Suppose $f\simeq\tilde{f}$ via operators $u:H\to\tilde{H}$ and $\tilde{u}:\tilde{H}\to H$, so $\boldsymbol{\rho}(u\circ r)\asymp\boldsymbol{\rho}(\tilde{r})$ (both defined on $\tilde{H}$) and $\boldsymbol{\rho}(\tilde{u}\circ\tilde{r})\asymp\boldsymbol{\rho}(r)$ (both defined on $H$). To show $\tilde{r}\simeq r$, take $v:=\tilde{u}:\tilde{H}\to H$ and $\tilde{v}:=u:H\to\tilde{H}$. \\
                  The first condition for $\tilde{r}\simeq r$ is $\boldsymbol{\rho}(v\circ\tilde{r})\asymp\boldsymbol{\rho}(r)$, i.e., $\boldsymbol{\rho}(\tilde{u}\circ\tilde{r})\asymp\boldsymbol{\rho}(r)$, which holds by assumption. \\
                  The second condition is $\boldsymbol{\rho}(\tilde{v}\circ r)\asymp\boldsymbol{\rho}(\tilde{r})$, i.e., $\boldsymbol{\rho}(u\circ r)\asymp\boldsymbol{\rho}(\tilde{r})$, which also holds by assumption. \\
                  Hence, $\tilde{r}\simeq r$.
            \item \textbf{Transitivity.} Suppose $f\simeq s$ via $u_1:H\to K$ and $\tilde{u}_1:K\to H$, with $s\in \mathsf{\mathcal{B}}(\Lambda,K)$, so $\boldsymbol{\rho}(u_1\circ r)\asymp\boldsymbol{\rho}(s)$ (both on $K$) with constants $K_1,K_2>0$:
                  \[ \forall z\in K:\quad K_1\boldsymbol{\rho}(s)(z)\leq\boldsymbol{\rho}(u_1\circ r)(z)\leq K_2\boldsymbol{\rho}(s)(z), \]
                  and $\boldsymbol{\rho}(\tilde{u}_1\circ s)\asymp\boldsymbol{\rho}(r)$ (both on $H$) with constants $\tilde{K}_1,\tilde{K}_2>0$:
                  \[ \forall z \in H:\quad \tilde{K_1} \boldsymbol{\rho}(r)(z)\leq\boldsymbol{\rho}(\tilde{u_1}\circ s)(z)\leq \tilde{K_2} \boldsymbol{\rho}(r)(z), \]
                  Further suppose $s\simeq\tilde{f}$ via $u_2:K\to\tilde{H}$ and $\tilde{u}_2:\tilde{H}\to K$, so $\boldsymbol{\rho}(u_2\circ s)\asymp\boldsymbol{\rho}(\tilde{r})$ (both on $\tilde{H}$) with constants $L_1,L_2>0$:
                  \[ \forall z \in\tilde{H}:\quad L_1\boldsymbol{\rho}(\tilde{r})(z)\leq\boldsymbol{\rho}(u_2\circ s)(z)\leq L_2\boldsymbol{\rho}(\tilde{r})(z), \]
                  and $\boldsymbol{\rho}(\tilde{u}_2\circ\tilde{r})\asymp\boldsymbol{\rho}(s)$ (both on $K$) with constants $\tilde{L}_1,\tilde{L}_2>0$:
                  \[ \forall z \in K:\quad \tilde{L}_1\boldsymbol{\rho}(s)(z)\leq\boldsymbol{\rho}(\tilde{u_2} \circ \tilde{r})(z)\leq \tilde{L}_2\boldsymbol{\rho}(s)(z), \]
                  To show $f\simeq\tilde{f}$, take $u:=u_2u_1:H\to\tilde{H}$ and $\tilde{u}:=\tilde{u}_1\tilde{u}_2:\tilde{H}\to H$. \\
                  Consider the first condition $\boldsymbol{\rho}(u\circ r)\asymp\boldsymbol{\rho}(\tilde{r})$ (both on $\tilde{H}$). \\
                  From $f\simeq s$, we have
                  \[ K_1\boldsymbol{\rho}(s)(z)\leq\boldsymbol{\rho}(u_1\circ r)(z)\leq K_2\boldsymbol{\rho}(s)(z), \]
                  Consequently
                  \[ \boldsymbol{\rho}(u\circ r)(z) = \boldsymbol{\rho}(r)(u^* z) = \boldsymbol{\rho}(r)(u_1^* u_2^* z) =  \boldsymbol{\rho}(u_1 \circ r) (u_2^* z) \]
                  so
                  \[ K_1\boldsymbol{\rho}(s)(u_2^*z)\leq\boldsymbol{\rho}(u\circ r)(z)\leq K_2\boldsymbol{\rho}(u_2\circ s)(z). \]
                  But $\boldsymbol{\rho}(s)(u_2^* z) = \boldsymbol{\rho}(u_2 \circ s)$, so
                  \[ K_1 \boldsymbol{\rho}(u_2 \circ s) \leq\boldsymbol{\rho}(u\circ r)(z)\leq K_2 \boldsymbol{\rho}(u_2 \circ s). \]
                  From $s\simeq\tilde{f}$, we have
                  $L_1\boldsymbol{\rho}(\tilde{r})(z)\leq\boldsymbol{\rho}(u_2\circ s)(z)\leq L_2\boldsymbol{\rho}(\tilde{r})(z)$, so
                  \[
                    K_1L_1\boldsymbol{\rho}(\tilde{r})(z)\leq\boldsymbol{\rho}(u\circ r)(z)\leq K_2L_2\boldsymbol{\rho}(\tilde{r})(z).
                  \]
                  Thus, $\boldsymbol{\rho}(u\circ r)\asymp\boldsymbol{\rho}(\tilde{r})$ with constants $K_1L_1,K_2L_2>0$. \\
                  The second condition $\boldsymbol{\rho}(\tilde{u}\circ\tilde{r})\asymp\boldsymbol{\rho}(r)$ (both on $H$) can be proven similarly. \\
                  Hence, $f\simeq\tilde{f}$.
          \end{itemize}
  \end{itemize}
\end{proof}

\mbox{} \\
Before moving on, let's prove the following.

\begin{definition}
  Let $M$ be a preordered monoid. \\
  We construct a 2-category $\mathsf{PreOrd\text{-}M^+Set}$, called the \textbf{2-category of preordered $M^+$-sets}, as follows.

  \begin{itemize}
    \item \textbf{0-cells.} preordered sets $X$ endowed with a monoid action by $M$ such that $(\forall m \in M^+)(\forall x \geq y) : m.x \geq m.y$. 
    \item \textbf{1-cells.} For each 0-cells $X_1, X_2$, 
    \begin{align*} \mathsf{Hom( X_1, X_2)} := &\{ \text{monotone maps  from } X_1 \text{ to } X_2 \text{ such that } \\
                               & (\forall m \in M^+)(\forall x \in X_1) :  f(m.x) = m.f(x) \}
    \end{align*}
    \item \textbf{2-cells.} For any two parallel morphisms $f,g \in Hom(X_1,X_2)$,
  \[ \mathsf{Hom(f,g)} = \{ c \in Z(M^+) \text{ such that } (\forall x,y \in X_1) : x \geq y \Rightarrow  f(x) \geq c g(y) \}. \]

  \item \textbf{Composition of 1-cells.}
          Composition of 1-morphisms is ordinary function composition, which is associative, unital, and preserves the monotonicity and equivariance of 1-morphisms.
    \item \textbf{Identity 1-cells.}
    For each object \(X\), the identity 1-morphism \(\operatorname{id}_X\) is the identity map, which is clearly monotone and equivariant.
    
    

            \item \textbf{Identity 2-cells.} \\
                  The identity 2-morphism \(1_\alpha\) on a 1-morphism \(f\) is the unit \(1 \in M^+\). \\
                  Indeed, since \(1 \in Z(M^+)\) and \(f(x) \geq 1.f(y) = f(y)\) by monotonicity of \(f\), we have \(1 \in \operatorname{Hom}(f, f)\), and 1 is the multiplicative unit of $M^+$.

            \item \textbf{Vertical composition of 2-cells.} \\
                  Given \(c \in \operatorname{Hom}(f, g)\) and \(d \in \operatorname{Hom}(g, h)\), the vertical composite is defined as \(d \circ c = cd \) (product in \(M^+\)).
                  \begin{itemize} 
                    \item \textbf{Well-definedness}:
                  \begin{itemize}
                    \item Since \(c, d \in Z(M^+)\), their product \(cd\) is also in \(Z(M^+)\).
                    \item For any \(x \geq y\), we have \(f(x) \geq c.g(y)\) and \(g(y) \geq d.h(y)\). Applying monotonicity of the action by \(c \in M^+\) to the second inequality gives \(c.g(y) \geq cd.h(y)\). By transitivity, \(f(x) \geq cd.h(y)\). Thus \(d \circ c = cd \in \operatorname{Hom}(f, h)\).
                  \end{itemize}

            \item \textbf{Associativity.} Follows from the monoid structure of \(M^+\).
          \item \textbf{Unit laws.} Follows from the monoid structure of \(M^+\).
          
                  \end{itemize}


            \item \textbf{Horizontal composition of 2-cells.} \\
                  Given \(c \in \operatorname{Hom}(f, g)\) (where \(f, g: X_1 \to X_2\)) and \(d \in \operatorname{Hom}(h, i)\) (where \(h, i: X_2 \to X_3\)), the horizontal composite is \(d * c = cd\).
                  \begin{itemize}
                  \item \textbf{Well-definedness}:
                  \begin{itemize}
                    \item  Since \(c, d \in Z(M^+)\), their product \(cd\) is also in \(Z(M^+)\).
                    \item We must check that \(cd \in \operatorname{Hom}(h \circ f, i \circ g)\).  \\
                          For any \(x \geq y\) in \(X_1\), we have \(f(x) \geq c.g(y)\). \\
                          Applying the monotone map \(h\) gives \(h(f(x)) \geq h(c.g(y))\). \\
                          By equivariance, \(h(c.g(y)) = c.h(g(y))\), so \(h(f(x)) \geq c.h(g(y))\).  \\
                          Since \(g(y) \geq g(y)\) (reflexivity), the condition on \(d\) gives \(h(g(y)) \geq d.i(g(y))\). Monotonicity of the action by \(c\) yields \(c.h(g(y)) \geq cd.i(g(y))\).  \\
                          By transitivity, \(h(f(x)) \geq cd.i(g(y))\). Thus \(d * c = cd \in  \operatorname{Hom}(h \circ f, i \circ g) \).
                  \end{itemize}

            \item \textbf{Associativity.} Follows from the monoid structure of \(M^+\).
            
            \item \textbf{Unit laws.} Follows from the monoid structure of \(M^+\).
          \end{itemize}

            \item \textbf{Interchange law.} \\
                  Given \(c \in \operatorname{Hom}(f, f')\), \(c' \in \operatorname{Hom}(f', f'')\), \(d \in \operatorname{Hom}(g, g')\), and \(d' \in \operatorname{Hom}(g', g'')\), we have
                  \[
                    (d' \circ d) * (c' \circ c) = (dd') * (cc') = cc'dd'
                  \]
                  and
                  \[
                    (d' * c') \circ (d * c) = (c'd') \circ (cd) = cdc'd'
                  \]
                  Since \(c', d \in Z(M^+)\), they commute, so \(cc'dd' = cdc'd'\). \\
                  Thus the two expressions are equal.

          \item \textbf{Associativity constraint.}

          For any triple of composable 1-cells
          $X_1 \to X_2 \to X_3 \to X_4$, the associator
          \[
            \alpha_{X_1,X_2,X_3,X_4}
          \]
          is the 2-cell $1 \in Z(M^+)$.
          The pentagon coherence condition holds trivially, since all 2-cells involved are equal to $1 \in Z(M^+)$.

    \item \textbf{Unit constraints.}
          Similarly, the left and right unitors are the 2-cell $1 \in Z(M^+)$, and the triangle coherence condition holds trivially.
  \end{itemize}
  All axioms of a 2-category are satisfied. Therefore, \(\mathsf{PreOrd\text{-}M^+Set}\) is a (strict) 2-category.  
\end{definition}

\mbox{} \\
Now, let's show that the equivalence relation of definition \ref{definition-operatorial-equivalence} fits into the general framework of definition \ref{definition-master}. \\
Let $\mathsf{Hilb_\mathbb{F}}$ be one of the categories $\mathsf{Hilb_\mathbb{F}^\text{inj}}$, or $\mathsf{Hilb_\mathbb{F}^\text{inj,cl}}$, or $\mathsf{Hilb_\mathbb{F}^\text{iso}}$, and $\mathsf{(PreOrd\text{-}\mathbb{R}^+_*Set)^{op}}$ be the opposite category of the 2-category $\mathsf{PreOrd\text{-}\mathbb{R}^+_*Set}$. \\
Define the function $\sigma : \mathsf{Hilb_\mathbb{F}} \to \mathsf{(PreOrd\text{-}\mathbb{R}^+_*Set)^{op}}$ as follows. \\
We set
\[ \sigma(H) = \{ \text{seminorms on } H \text{ dominated by } \lVert \cdot \rVert_H \} \]
for any Hilbert space $H$ over $\mathbb{F}$. \\
Moreover, we set
\[ \sigma(m) = \begin{cases} \{ \text{seminorms on } H_2 \text{ dominated by } \lVert \cdot \rVert_{H_2} \} &\to \{ \text{seminorms on } H_1 \text{ dominated by } \lVert \cdot \rVert_{H_1} \} \\ \quad \quad \quad \quad \quad \quad \phi &\mapsto \left( \underset{\lVert x \rVert_{H_2} = 1}{\operatorname{sup}} \phi(x) \right) \lVert m(\cdot) \rVert_{H_2} \end{cases} \]
for any continuous operator $m : H_1 \to H_2$. \\
It is clear that $\sigma(m)$ is well-defined, $\mathbb{R}^+_*$-equivariant, and monotonic, so $\sigma(m)$ is a 1-morphism in $\mathsf{(PreOrd\text{-}\mathbb{R}^+_*Set)^{op}}$. \\
Notice that $\sigma$ is only a function and not a functor, since for any $m : H_1 \to H_2$, $\overline{m} : H_2 \to H_3$, and seminorm $\phi$ on $H_3$ dominated by $\lVert \cdot \rVert_{H_3}$
\[ \sigma(m) \circ \sigma(\overline{m}) (\phi) =  \left( \underset{\lVert x \rVert_{H_2} = 1}{\operatorname{sup}} \sigma(\overline{m})(\phi)(x) \right) \lVert m(\cdot) \rVert_{H_2} =  \left( \underset{\lVert x \rVert_{H_3} = 1}{\operatorname{sup}} \phi(x) \right) \lVert \overline{m} \rVert_{H_2;H_3} \lVert m(\cdot) \rVert_{H_2} \]
is not necessarily equal to
\[ \sigma(\overline{m} \circ m)(\phi)  =  \left( \underset{\lVert x \rVert_{H_3} = 1}{\operatorname{sup}} \phi(x) \right) \lVert (\overline{m} \circ m)(\cdot) \rVert_{H_2} \]
What's more, define the functor $\tau_1 : \mathsf{Hilb_\mathbb{F}} \to \mathsf{(PreOrd\text{-}\mathbb{R}^+_*Set)^{op}}$ as follows. \\
We set
\[ \tau_1(H) = \{ \text{seminorms on } H \text{ dominated by } \lVert \cdot \rVert_H \} \]
for any Hilbert space $H$ over $\mathbb{F}$. \\
Moreover, we set
\[ \tau_1(m) = \begin{cases} \{ \text{seminorms on } H_2 \text{ dominated by } \lVert \cdot \rVert_{H_2} \} &\to \{ \text{seminorms on } H_1 \text{ dominated by } \lVert \cdot \rVert_{H_1} \} \\ \quad \quad \quad \quad \quad \quad \quad \phi &\mapsto \quad \quad \quad \quad \phi \circ m \end{cases} \]
for any continuous operator $m : H_1 \to H_2$. \\
It is clear that $\tau_1(m)$ is well-defined, $\mathbb{R}^+_*$-equivariant, and monotonic, so $\tau_1(m)$ is a 1-morphism in $\mathsf{(PreOrd\text{-}\mathbb{R}^+_*Set)^{op}}$. \\
Moreover, $\tau_1$ is a functor from $\mathsf{Hilb_\mathbb{F}}$ to $\mathsf{(PreOrd\text{-}\mathbb{R}^+_*Set)^{op}}$, since for any $m : H_1 \to H_2$, $\overline{m} : H_2 \to H_3$, and seminorm $\phi$ on $H_3$ dominated by $\lVert \cdot \rVert_{H_3}$
\[ \tau_1(\overline{m})\tau_1(m)(\phi) = \tau_1(m)(\phi) \circ \overline{m} = \phi \circ m \circ \overline{m} = \tau_1(m \circ \overline{m})(\phi). \]
Finally, define the functor $\tau_2 : \mathsf{Hilb_\mathbb{F}} \to \mathsf{(PreOrd\text{-}\mathbb{R}^+_*Set)^{op}}$ as follows.
We set
\[ \tau_2(H) = \{ \text{seminorms on } H \text{ dominated by } \lVert \cdot \rVert_H \} \]
for any Hilbert space $H$ over $\mathbb{F}$. \\
Moreover, we set
\[ \tau_2(m) = \begin{cases} \{ \text{seminorms on } H_2 \text{ dominated by } \lVert \cdot \rVert_{H_2} \} &\to \{ \text{seminorms on } H_1 \text{ dominated by } \lVert \cdot \rVert_{H_1} \} \\ \quad \quad \quad \quad \quad \quad \quad \phi &\mapsto \quad \quad \quad \quad \left( \underset{\lVert x \rVert_{H_2} = 1}{\operatorname{sup}} \phi(x) \right) \lVert \cdot \rVert_{H_1} \end{cases} \]
for any continuous operator $m : H_1 \to H_2$. \\
It is clear that $\tau_2(m)$ is well-defined, $\mathbb{R}^+_*$-equivariant, and monotonic, so $\tau_2(m)$ is a 1-morphism in $\mathsf{(PreOrd\text{-}\mathbb{R}^+_*Set)^{op}}$. \\
Moreover, $\tau_2$ is a functor from $\mathsf{Hilb_\mathbb{F}}$ to $\mathsf{(PreOrd\text{-}\mathbb{R}^+_*Set)^{op}}$, since for any $m : H_1 \to H_2$, $\overline{m} : H_2 \to H_3$, and seminorm $\phi$ on $H_3$ dominated by $\lVert \cdot \rVert_{H_3}$
\[ \tau_2(m)(\phi)(h_1) = \left( \underset{\lVert x \rVert_{H_2} = 1}{\operatorname{sup}} \phi(x) \right) \lVert h_1 \rVert \]
and so
\begin{align*}
  \tau_2(\overline{m}) \tau_2(m) (\phi) & = \left( \underset{\lVert x \rVert_{H_2} = 1}{\operatorname{sup}} (\tau_2(m) \phi)(x) \right)  \lVert \cdot \rVert_{H_1}                                                                                           \\
                                        & =  \left( \underset{\lVert x \rVert_{H_2} = 1}{\operatorname{sup}} \left( \underset{\lVert y \rVert_{H_3} = 1}{\operatorname{sup}} \phi(y) \right)  \lVert x \rVert_{L^2(H_2)} \right)  \lVert  \cdot \rVert_{H_1} \\
                                        & = \left( \underset{\lVert y \rVert_{H_3} = 1}{\operatorname{sup}} \phi(y) \right)  \lVert \cdot \rVert_{H_1}                                                                                                       \\
                                        & = \tau_2(m \circ \overline{m}) (\phi)
\end{align*}

\begin{proposition}
  Two functions $f \in \mathsf{\mathcal{B}}(\Omega,H)$ and $\tilde{f} \in \mathsf{\mathcal{B}}(\tilde{\Omega},\tilde{H})$ are equivalent as in definition \ref{definition-operatorial-equivalence} iff
  \[ T_f : H \to L^2(\Omega) \]
  and
  \[ T_{\tilde{f}} : \tilde{H} \to L^2(\tilde{\Omega}) \]
  are equivalent as in definition \ref{definition-master}, with respect to the present data $(\mathsf{Hilb_\mathbb{F}},\mathsf{(PreOrd\text{-}\mathbb{R}^+_*Set)^{op}},\sigma,\tau_1,\tau_2)$.
\end{proposition}

\begin{proof}
  First, for any seminorm $\phi$ on $L^2(\tilde{\Omega})$ dominated by $\lVert \cdot \rVert_{L^2(\tilde{\Omega})}$, we have
  \[ \sigma(T_{\tilde{f}})(\phi) =  \left( \underset{\lVert x \rVert_{L^2(\tilde{\Omega})} = 1}{\operatorname{sup}} \phi(x) \right) \lVert T_{\tilde{f}}(\cdot) \rVert_{L^2(\tilde{\Omega})} = \left( \underset{\lVert x \rVert_{L^2(\tilde{\Omega})} = 1}{\operatorname{sup}} \phi(x) \right) \rho(T_{\tilde{f}}). \]
  Moreover, for all $u_1 : \tilde{H} \to H$, $u_2 : L^2(\Omega) \to L^2(\tilde{\Omega})$, and any seminorm $\phi$ on $L^2(\tilde{\Omega})$ dominated by $\lVert \cdot \rVert_{L^2(\tilde{\Omega})}$, we have
  \begin{align*}
    \left[ \tau_1(u_1) \circ \sigma(T_f) \circ \tau_2(u_2) \right] (\phi) & = \left[ \sigma(T_f)(\tau_2(u_2)(\phi)) \right] \circ u_1 \\ & =  \left( \underset{\lVert x \rVert_{L^2(\Omega)} = 1}{\operatorname{sup}}  (\tau_2(u_2)(\phi)(x) \right) \lVert T_{f}(\cdot) \rVert_{L^2(\Omega)} \circ u_1 & = \left( \underset{\lVert x \rVert_{L^2(\Omega)} = 1}{\operatorname{sup}}  \left( \underset{\lVert y \rVert_{L^2(\tilde{\Omega})} = 1}{\operatorname{sup}} \phi(y) \right) \lVert x \rVert_{L^2(\Omega)} \right) \lVert (T_{f} \circ u_1)(\cdot) \rVert_{L^2(\Omega)} \\ & = \left( \underset{\lVert y \rVert_{L^2(\tilde{\Omega})} = 1}{\operatorname{sup}} \phi(y) \right) \rho(T_f \circ u_1).
  \end{align*}
  Now,
  \begin{itemize}
    \item If there exist $u_1 : \tilde{H} \to H$, $u_2 : L^2(\Omega) \to L^2(\tilde{\Omega})$, and two 2-morphisms
          \[ c : \tau_1(u_1) \circ \sigma(T_f) \circ \tau_2(u_2) \Rightarrow \sigma(T_{\tilde{f}}) \]
          and
          \[ \tilde{c} : \sigma(T_{\tilde{f}})  \Rightarrow \tau_1(u_1) \circ \sigma(T_f) \circ \tau_2(u_2), \]
          then evaluating at $\phi := \lVert \cdot \rVert_{L^2(\tilde{\Omega})}$, we obtain
          \[ \rho(T_f \circ  u_1) = \tau_1(u_1) \circ \sigma(T_f) \circ \tau_2(u_2) (\phi) \geq c  \sigma(T_{\tilde{f}})(\phi) = c \rho(T_{\tilde{f}}) \]
          and so
          \[ \frac{1}{c} \rho(T_f \circ  u_1) \geq  \rho(T_{\tilde{f}}), \]
          and
          \[
            \rho(T_{\tilde{f}}) = \sigma(T_{\tilde{f}})(\phi) \geq c' \tau_1(u_1) \circ \sigma(T_f) \circ \tau_2(u_2) (\phi) = c' \rho(T_f \circ  u_1).
          \]
          Thus $\rho(T_f \circ  u_1) \asymp \rho(T_{\tilde{f}})$.

    \item Conversely, if $\rho(T_f \circ  u_1) \asymp \rho(T_{\tilde{f}})$, then for any two seminorms $\phi,\tilde{\phi}$ on $L_2(\tilde{\Omega})$ dominated by $\lVert \cdot \rVert_{L_2(\tilde{\Omega})}$ and such that $\tilde{\phi} \geq \phi$, we have
          \begin{align*}
            \tau_1(u_1) \circ \sigma(T_f) \circ \tau_2(u_2) (\tilde{\phi}) & = \left( \underset{\lVert y \rVert_{L^2(\tilde{\Omega})} = 1}{\operatorname{sup}} \tilde{\phi}(y) \right) \rho(T_f \circ u_1) \\
                                                                           & \geq k_2  \left( \underset{\lVert y \rVert_{L^2(\tilde{\Omega})} = 1}{\operatorname{sup}} \phi(y) \right) \rho(T_{\tilde{f}}) \\
                                                                           & = k_2  \sigma(T_{\tilde{f}})(\phi)
          \end{align*}
          and
          \begin{align*}
            \sigma(T_{\tilde{f}})(\tilde{\phi}) & = \left( \underset{\lVert x \rVert_{L^2(\Omega)} = 1}{\operatorname{sup}}  \tilde{\phi}(x) \right) \rho(T_{\tilde{f}}) \\
                                                & \geq k_1 \left( \underset{\lVert x \rVert_{L^2(\Omega)} = 1}{\operatorname{sup}}  \phi(x) \right) \rho(T_f \circ  u_1) \\
                                                & = k_1 \tau_1(u_1) \circ \sigma(T_f) \circ \tau_2(u_2) (\phi)
          \end{align*}
          for somes positive constants $k_1,k_2 > 0$ whose existence is guaranteed by $\rho(T_f \circ  u_1) \asymp \rho(T_{\tilde{f}})$.
  \end{itemize}
  Similarly, the existence of $\tilde{u}_1 :  H \to \tilde{H}$, $\tilde{u}_2 : L^2(\tilde{\Omega}) \to L^2(\Omega)$, and two 2-morphisms
  \[ d : \tau_1(\tilde{u}_1) \circ \sigma(T_{\tilde{f}}) \circ \tau_2(\tilde{u}_2) \Rightarrow \sigma(T_f) \]
  and
  \[ \tilde{d} : \sigma(T_f)  \Rightarrow \tau_1(\tilde{u}_1) \circ \sigma(T_{\tilde{f}}) \circ \tau_2(\tilde{u}_2) \]
  is equivalent to
  \[ \rho(T_{\tilde{f}} \circ u_2) \asymp \rho(T_f). \]
\end{proof}

Finally, two questions deserve to be asked.

\begin{problem}
Characterize the equivalence classes of the equivalence relation $\simeq$ of definition \ref{definition-operatorial-equivalence}. There is a natural equivalence class consisting of frames (or equivalently, orthonormal bases), but what about the other equivalence classes? 
\end{problem}

\begin{problem}
Generalize the equivalence relation $\simeq$ of definition \ref{definition-operatorial-equivalence} to the Banach space setting. Here, the difficulty is that continuous Bessel families and their analysis operators depend on extra structure. 
\end{problem}



\section{Conclusion}
\label{section-conclusion}

Our generalized equivalence relation (definition \ref{definition-master}) provides a flexible and conceptually clean framework that extends many classical equivalence relations in mathematics (see definitions \ref{definition-group-equivalence}, \ref{definition-pullback-equivalence-relation}, and \ref{definition-operatorial-equivalence}). \\
More broadly, this approach provides a unifying language for equivalence problems where the relevant structure is encoded by 2-categories and observables rather than explicit symmetries. \\ \\

\section*{Ethics declaration}
Not applicable.

\section*{Funding declaration}
This research received no specific grant from any funding agency in the public, commercial, or not-for-profit sectors.

\section*{Competing interests declaration}
The author declares no competing interests.

%%%%%%%%%%%%%%%%%%%%
\nocite{*} % include non-cited references
\bibliographystyle{amsplain}
\bibliography{references}
%%%%%%%%%%%%%%%%%%%%

\Addresses

\end{document}
